\section{Completion of the proof}\label{sec:assembly}

We now combine the estimates, keeping the order of constants explicit.
This first proves Theorem~\ref{thm:survivors}; the deduction of
Theorem~\ref{thm:main} then uses only an upper sieve.

\subsection{The error budget}

\begin{proof}[Proof of Theorem~\ref{thm:survivors}]
By Proposition~\ref{prop:reference-margin}, there is an absolute
$c_L>0$ such that
\begin{equation}\label{eq:assembly-reference}
 P_\ee(r_0,B)\ge\frac{c_L}{B^2}
\end{equation}
for all sufficiently large $z$. Let $C_{\rm abs}$ be an absolute
constant such that
\[
 \left|N_d/\mu_d-1\right|\le C_{\rm abs}
\]
at every included node; this follows from
\eqref{eq:tree-small-counts}.

For nodes with large gap, the relative error in
\eqref{eq:tree-small-counts} is $O(e^{-cr_d})$ after a fixed lower
threshold on $r_d$. The tail estimate in
Proposition~\ref{prop:prime-occupation} permits us to choose a fixed
$K$ so large that
\begin{equation}\label{eq:assembly-tail}
 \sum_{\substack{d\text{ included}\\r_d>K}}
 |N_d-\mu_d|
 \le\frac{c_L}{8}\frac{YV_0}{B^2}
\end{equation}
eventually. Standard paths contain the included paths, so the
nonnegative tail estimate applies. This choice depends only on the
fixed reference margin and the absolute small-prime error constants.

The same proposition gives a constant $C_K$ for the total harmonic
mass of nodes with $r_d\le K$. Choose $\eps>0$ such that
$\eps C_K<c_L/8$. The compact nodes with relative error at most
$\eps$ then contribute at most
\begin{equation}\label{eq:assembly-good-compact}
 \sum_{\substack{d\text{ included}\\r_d\le K\\
                    |N_d/\mu_d-1|\le\eps}}
 |N_d-\mu_d|
 \le\frac{c_L}{8}\frac{YV_0}{B^2}.
\end{equation}

It remains to deal with compact nodes having relative error greater
than $\eps$. Fix a target $\beta>0$ so small that
\begin{equation}\label{eq:assembly-beta}
 C_{\rm abs}\beta<\frac{c_L}{8}.
\end{equation}
We claim that the stopped tree can be chosen so that the retained
bad compact nodes have total harmonic mass at most $\beta/B^2$.
The following choices give this claim without requiring any constant
to be chosen again at an earlier stage.

First use Lemma~\ref{lem:edge-witness} to fix a larger threshold
$K_*$, a positive edge-discrepancy threshold, and a fixed maximum
number $H_e$ of candidate positions. Choose the adjustable regularity
exceptions in Proposition~\ref{prop:regular-boxes} small enough to
cost less than a prescribed small fraction of $\beta/B^2$.
This fixes the path-length bound and the isolated-prime interval.
For the moment leave the exception for near-threshold even prefixes
to the eventual choice of bin width.

The enlarged regular boxes have total harmonic mass at most
$C'_K/B^2$, with $C'_K$ independent of sufficiently small fixed
$\xi$. Choose the exception fraction $\sigma$ in
Proposition~\ref{prop:inverse} sufficiently small that
\[
 H_eC'_K\sigma
\]
is less than another prescribed small fraction of $\beta$.
The inverse proposition now fixes a positive alignment fraction and
a list bound $F$, independently of the eventual sufficiently small
fixed width. Each list is attached to its entire candidate box.

Next choose the variance exception fraction in
Proposition~\ref{prop:structured-variance} small enough relative to
$H_e$, $C'_K$, $F$ and $\beta$. A prime-count fraction in the isolated
bin converts to a harmonic fraction at a bounded absolute cost.
After this conversion and the finite unions over candidates and
list elements, the total variance exception costs another prescribed
small fraction of $\beta/B^2$. The variance proposition fixes a
constant $C_s$. Corollary~\ref{cor:hard-endpoints} then shows that
all remaining bad compact endpoints have a hard rational satisfying
the global eligibility bounds and alignment of every factor with
exponent greater than $C_s$.

Choose the absolute owner threshold $\eta$ required by
Proposition~\ref{prop:stopped-comparison}. Choose $b_0$ sufficiently
large in terms of $C_s$ and $K$ and for the local reference and prime
estimates in that proposition. Finally choose $b_1/b_0$ sufficiently
large that Proposition~\ref{prop:hard-stopping} bounds the hard
endpoints that remain outside stopped subtrees by the remaining
allocated small fraction of $\beta/B^2$. This uses the marked-visit
estimate with a prescribed failure probability, as well as the
vanishing owner and wrong-tag exceptions for the fixed candidate
lists.

All constants except the bin width have now been fixed. Choose a
fixed $\xi>0$ sufficiently small for the interval-freezing errors,
candidate-size slack, representative stopping margins and all
earlier small-width requirements. In particular, the near-threshold
exception, of cost $O(\Delta)+o(1)$ in the compact-history
probability measure with $\Delta=2C_{\rm len}\xi$, can be made as
small as the error budget requires. Then take $z$ sufficiently large
for every preceding fixed choice. Asymptotic bounds whose thresholds
depend on $\xi$ are used only at this final stage.

The sum of all these finitely many allocated costs is at most
$\beta/B^2$, proving the claim. More explicitly, after removing the
regularity exceptions, every retained bad endpoint has a witnessing
edge. Outside the inverse exceptions it has an alignment in the
fixed list. Outside the variance exceptions it is hard. Outside
the hard-stopping exceptions it lies in a stopped subtree and is
therefore not retained. This exhausts the cases. Fractional
exceptions are summed using $H_eC'_K/B^2$, not the number of boxes.

The claim and \eqref{eq:assembly-beta} show that retained bad compact
nodes contribute at most $(c_L/8)YV_0/B^2$ in absolute error.
Together with \eqref{eq:assembly-tail} and
\eqref{eq:assembly-good-compact}, this gives
\begin{equation}\label{eq:assembly-retained}
 \sum_{d\in\mathcal T}|N_d-\mu_d|
 \le\frac{3c_L}{8}\frac{YV_0}{B^2}
\end{equation}
for the retained expanded nodes of the chosen stopped tree.
Proposition~\ref{prop:stopped-comparison} gives
\[
 \sum_{d\in\mathcal S}
 \bigl(S_d(b_d)-\mu_dP_\ee(r_d,b_d)\bigr)\ge0.
\]
Apply \eqref{eq:stopped-comparison-identity},
Lemma~\ref{lem:tree-lower-bound} and
\eqref{eq:assembly-reference}. We obtain
\[
 S_1(B)\ge\left(1-\frac38\right)c_L\frac{YV_0}{B^2}
 \ge\frac{c_L}{2}\frac{YV_0}{B^2}.
\]
Every constant is independent of the prescribed residue classes.
This proves Theorem~\ref{thm:survivors}.
\end{proof}

\subsection{Removing the remaining divisor primes}

The factor $\log w$ in the survivor estimate allows a cutoff smaller
than $k\log k$.  We choose it so that the loss from the at most $k$
remaining divisor primes is a fixed fraction of the surviving count.

\begin{proof}[Proof of Theorem~\ref{thm:main}]
By Theorem~\ref{thm:survivors} and Mertens' formula, there is an
absolute $c_0>0$ such that, for every sufficiently large real $z$,
every assignment of one forbidden class at each prime $p\le z$
leaves at least
\begin{equation}\label{eq:assembly-survivor-lower}
 c_0\frac{Y\log w}{L^2}
\end{equation}
integers in $[1,Y]$.  Indeed,
$V_0/B^2\sim e^{-\gamma}\log w/L^2$.

Let $k$ be sufficiently large, put
$\ell=\log k$ and $M=\log\log k$, and take
\begin{equation}\label{eq:assembly-jacobsthal-cutoff}
 z=A_0\frac{k\log k}{\log\log k},
\end{equation}
where the absolute constant $A_0>1$ will be fixed below, before $k$
tends to infinity.  With the parameters of \eqref{eq:parameters},
\begin{equation}\label{eq:assembly-jacobsthal-scales}
 L\sim\ell,\qquad \log w\sim M,\qquad
 Y\sim A_0^2\frac{k^2}{M^2},\qquad
 X:=\frac Yz\sim A_0\frac{k}{\ell M},\qquad
 \log X\sim\ell.
\end{equation}
In particular $z\to\infty$ and $X\to\infty$ for every such fixed
$A_0$.

Let $n$ be any positive integer with $\omega(n)\le k$, and let the
given interval of $Y$ consecutive integers begin at an arbitrary
integer $a$.  Represent its integers as $a+j-1$, $1\le j\le Y$.
For each prime $p\mid n$ with $p\le z$, prescribe the forbidden class
\[
 j\equiv1-a\pmod p.
\]
At the other primes $p\le z$, prescribe any class.  The estimate
\eqref{eq:assembly-survivor-lower} counts integers that avoid every
divisor prime at most $z$, together with these additional restrictions.

Fix a prime divisor $q>z$ of $n$.  Its divisibility class has at most
$1+Y/q\le1+X$ representatives in the coordinate interval.  Enclose
them in a progression segment whose coordinate length is
$\lceil X\rceil+2$.  Fix a sufficiently small absolute $\theta>0$
and sieve this segment by primes at most $X^\theta$.  This cutoff is
below $z$ for large $k$, the step $q$ is invertible modulo every sieve
prime, and the integers already counted in
\eqref{eq:assembly-survivor-lower} avoid its prescribed class there.
Lemma~\ref{lem:small-sieve} and Mertens' formula therefore give
\begin{equation}\label{eq:large-prime-deletion}
 \#\{\text{these survivors divisible by }q\}
 \le C_1\left(1+\frac X{\log X}\right)
\end{equation}
for an absolute $C_1$, uniform in $q$, $A_0$, the classes, and the
translation once the displayed length conditions hold.  More
explicitly, the enclosing coordinate length is comparable to $X$
and at least $(X^\theta)^{1/(2\theta)}$ for large $X$, so the fixed
exponent upper bound in Lemma~\ref{lem:small-sieve} applies.
Enclosing a shorter or empty progression only increases the count;
in particular the argument also covers $q>Y$.

There are at most $k$ such further divisor primes.  For each fixed
$A_0$, \eqref{eq:assembly-jacobsthal-scales} gives
\begin{equation}\label{eq:all-large-deletions}
 \frac{k(1+X/\log X)}{Y\log w/L^2}
 =\frac1{A_0}+o(1).
\end{equation}
For the second term in the numerator, the quotient is
$kL^2/(z\log X\log w)\to1/A_0$; the term $k$ is negligible since
$Y\log w/L^2\asymp_{A_0}k^2/(\ell^2M)$.
Choose once and for all $A_0>2C_1/c_0$.  Equations
\eqref{eq:large-prime-deletion} and \eqref{eq:all-large-deletions}
then show that fewer than the number in
\eqref{eq:assembly-survivor-lower} are removed, for all sufficiently
large $k$.  At least one integer in the original interval is therefore
coprime to $n$.

This argument is uniform in the prime divisors of $n$ and in every
integer $a$, including negative starting points and intervals containing
zero.  Prime multiplicities do not enter.  The length estimate in
\eqref{eq:assembly-jacobsthal-scales}, with $A_0$ now absolute, gives
\[
 h(k)\ll\frac{k^2}{(\log\log k)^2}
       \ll\frac{k^2}{(\log\log(3k))^2}
\]
for all sufficiently large $k$, since
$\log\log(3k)/\log\log k\to1$.

To include the remaining finitely many values of $k$, suppose $n$ has
$r\le k$ distinct prime divisors.  Inclusion--exclusion gives at least
\[
 m\prod_{p\mid n}\left(1-\frac1p\right)-2^r
 \ge\frac m{k+1}-2^k
\]
coprime integers in every interval of $m$ consecutive integers.
For the product bound, order the distinct primes as
$p_1<\cdots<p_r$ and use $p_j\ge j+1$ to obtain
$\prod_{j=1}^r(1-1/p_j)\ge1/(r+1)$.  This also covers $r=0$
with the empty product equal to one.  The integer
$m_k=(k+1)2^k+1$ consequently bounds $h(k)$ for each remaining $k$.
The maximum of
\[
 \frac{m_k(\log\log(3k))^2}{k^2}
\]
over this finite set, together with the constant already obtained for
large $k$, supplies one absolute constant $C$ such that
$h(k)\le Ck^2/(\log\log(3k))^2$ for every integer $k\ge1$.
\end{proof}
